\documentclass{article}
\usepackage{tikz}
\usetikzlibrary{calc,angles,quotes,intersections,decorations.markings}
\def\R{3}
\DeclareMathOperator{\longueur}{longueur}
\tikzset{netchinh/.style={line width=.8pt}}
\begin{document}
\begin{bt}[Câu 3]
	Cho tam giác $ABC$ nội tiếp đường tròn $(O)$, $AH$ là đường cao.
	\begin{center}
		\begin {tikzpicture}[scale=.8,font=\footnotesize,line join=round,line cap=round,>=stealth]
			\def\gocA{110}
			\path (0,0) coordinate (O)
				(\gocA:\R) coordinate (A)
				(210:\R) coordinate (B)
				(-30:\R) coordinate (C)
				($(B)!(A)!(C)$) coordinate (H);
			\draw[netchinh] (A)--(B)--(C)--cycle (O) circle(\R) (A)--(H);
			\foreach \x/\y/\z in {A/H/C}
				\path pic[draw=black,angle radius=6pt] {right angle=\x--\y--\z};
			\foreach \d/\g in {O/-45,A/100,B/-135,C/-30,H/-90}
				\path[draw,fill=white] (\d) circle(1pt) +(\g:9pt) node {$\d$};
		\end {tikzpicture}
	\end{center}
	\loigiai{
		Hình minh họa giữ nguyên quan hệ $AH\perp BC$.
	}
\end{bt}
\end{document}
